Une sphère métallique de \qty{\v{r}}{cm} de rayon est suspendue dans une enceinte sous vide dont les parois sont très froides (on néglige le rayonnement qu'elles émettent). Un élément chauffant interne fournit \qty[round-mode=figures,round-precision=3]{\fpeval{\v{em}*5.67e-8*4*pi*\vsi{r}^2*\v{T}^4}}{W} à la sphère, dont la température se stabilise à \qty{\v{T}}{K}.

\begin{enumerate}[label=\alph*)]
    \item Expliquez pourquoi, à l'équilibre, la puissance de chauffage est égale à la puissance rayonnée.
    \item Déterminez l'émissivité de la surface de la sphère.
    \item Quelle puissance de chauffage faudrait-il pour que la sphère se stabilise à \qty{\fpeval{2*\v{T}}}{K} ?
\end{enumerate}
